% !TeX program = lualatex
% !TeX root = pred3.tex
% ============================================================
% Zajednički preamble za sva predavanja RWA kolegija
% Uključiti s: % ============================================================
% Zajednički preamble za sva predavanja RWA kolegija
% Uključiti s: % ============================================================
% Zajednički preamble za sva predavanja RWA kolegija
% Uključiti s: \input{../common/preamble.tex}
% ============================================================

\documentclass[aspectratio=149,11pt]{beamer}
\usetheme{CambridgeUS}

\usepackage{fontspec}              % radi samo sa XeLaTeX/LuaLaTeX
\setsansfont{Latin Modern Sans}    % sans-serif font
\setmonofont{Latin Modern Mono}    % monospace font za kod
\usepackage[croatian]{babel}
\usepackage{csquotes}              % \enquote{} → hrvatski navodnici „..."
\usepackage{lmodern}
\usepackage{booktabs}
\usepackage{listings}
\usepackage{xcolor}
\usepackage{hyperref}

% ---------- Autor i logo ----------
\author{mag.ing. Josip Torić (SIT UNIZD)}
\titlegraphic{\includegraphics[height=2.5cm]{../common/SIT_logo_tekst1.pdf}}

% ---------- Listings ----------
\definecolor{codebg}{RGB}{248,248,248}
\definecolor{codeframeclr}{RGB}{220,220,220}
\lstdefinelanguage{bash}{
  morekeywords={sudo,apt,brew,pip,pipx,uv,winget,python,git,flask,source,deactivate,code,gunicorn,docker,compose,heroku,render,railway,node,npm,pnpm,yarn,uvicorn,alembic,pytest,ruff,black,precommit},
  sensitive=false,
  morecomment=[l]{\#}
}
\lstdefinelanguage{python}{
  morekeywords={from,import,def,class,return,if,elif,else,for,in,while,try,except,raise,with,as,not,and,or,is,lambda,yield,True,False,None,await,async},
  sensitive=true,
  morecomment=[l]{\#},
  morestring=[b]',
  morestring=[b]"
}
\lstdefinelanguage{sql}{
  morekeywords={SELECT,FROM,WHERE,INSERT,INTO,VALUES,UPDATE,SET,DELETE,CREATE,TABLE,PRIMARY,KEY,FOREIGN,REFERENCES,UNIQUE,NOT,NULL,INDEX,ON,ALTER,DROP,CONSTRAINT},
  sensitive=false,
  morecomment=[l]{--}
}
\lstdefinelanguage{json}{
  morecomment=[l]{//},
  morestring=[b]",
  showstringspaces=false
}
\lstdefinelanguage{typescript}{
  morekeywords={import,export,from,const,let,var,function,return,if,else,for,of,in,while,switch,case,break,async,await,try,catch,finally,throw,new,class,this,extends,implements,interface,type,enum,readonly,public,private,protected,static,get,set,true,false,null,undefined,void,string,number,boolean,any,unknown,never,typeof,keyof,as},
  sensitive=true,
  morecomment=[l]{//},
  morecomment=[s]{/*}{*/},
  morestring=[b]',
  morestring=[b]",
  morestring=[b]`
}
\lstalias{javascript}{typescript}
\lstset{
  backgroundcolor=\color{codebg},
  basicstyle=\ttfamily\footnotesize,
  frame=single,
  rulecolor=\color{codeframeclr},
  frameround=tttt,
  breaklines=true,
  tabsize=2,
  showstringspaces=false,
  upquote=true,
  columns=fullflexible
}

% ---------- Helper macros ----------
\newcommand{\codefile}[1]{\texttt{#1}}
\newcommand{\kbd}[1]{\texttt{#1}}

% ---------- Code frame environment ----------
% Automatski dodaje fragile i shrink=5, koristi se umjesto \begin{frame}
% Primjer: \begin{codeframe}{Naslov} ... \end{codeframe}
% Primjer s opcijama: \begin{codeframe}[shrink=15]{Naslov} ... \end{codeframe}
\newenvironment{codeframe}[2][]
  {\begin{frame}[fragile,shrink=5,#1]{#2}}
  {\end{frame}}

% ============================================================

\documentclass[aspectratio=149,11pt]{beamer}
\usetheme{CambridgeUS}

\usepackage{fontspec}              % radi samo sa XeLaTeX/LuaLaTeX
\setsansfont{Latin Modern Sans}    % sans-serif font
\setmonofont{Latin Modern Mono}    % monospace font za kod
\usepackage[croatian]{babel}
\usepackage{csquotes}              % \enquote{} → hrvatski navodnici „..."
\usepackage{lmodern}
\usepackage{booktabs}
\usepackage{listings}
\usepackage{xcolor}
\usepackage{hyperref}

% ---------- Autor i logo ----------
\author{mag.ing. Josip Torić (SIT UNIZD)}
\titlegraphic{\includegraphics[height=2.5cm]{../common/SIT_logo_tekst1.pdf}}

% ---------- Listings ----------
\definecolor{codebg}{RGB}{248,248,248}
\definecolor{codeframeclr}{RGB}{220,220,220}
\lstdefinelanguage{bash}{
  morekeywords={sudo,apt,brew,pip,pipx,uv,winget,python,git,flask,source,deactivate,code,gunicorn,docker,compose,heroku,render,railway,node,npm,pnpm,yarn,uvicorn,alembic,pytest,ruff,black,precommit},
  sensitive=false,
  morecomment=[l]{\#}
}
\lstdefinelanguage{python}{
  morekeywords={from,import,def,class,return,if,elif,else,for,in,while,try,except,raise,with,as,not,and,or,is,lambda,yield,True,False,None,await,async},
  sensitive=true,
  morecomment=[l]{\#},
  morestring=[b]',
  morestring=[b]"
}
\lstdefinelanguage{sql}{
  morekeywords={SELECT,FROM,WHERE,INSERT,INTO,VALUES,UPDATE,SET,DELETE,CREATE,TABLE,PRIMARY,KEY,FOREIGN,REFERENCES,UNIQUE,NOT,NULL,INDEX,ON,ALTER,DROP,CONSTRAINT},
  sensitive=false,
  morecomment=[l]{--}
}
\lstdefinelanguage{json}{
  morecomment=[l]{//},
  morestring=[b]",
  showstringspaces=false
}
\lstdefinelanguage{typescript}{
  morekeywords={import,export,from,const,let,var,function,return,if,else,for,of,in,while,switch,case,break,async,await,try,catch,finally,throw,new,class,this,extends,implements,interface,type,enum,readonly,public,private,protected,static,get,set,true,false,null,undefined,void,string,number,boolean,any,unknown,never,typeof,keyof,as},
  sensitive=true,
  morecomment=[l]{//},
  morecomment=[s]{/*}{*/},
  morestring=[b]',
  morestring=[b]",
  morestring=[b]`
}
\lstalias{javascript}{typescript}
\lstset{
  backgroundcolor=\color{codebg},
  basicstyle=\ttfamily\footnotesize,
  frame=single,
  rulecolor=\color{codeframeclr},
  frameround=tttt,
  breaklines=true,
  tabsize=2,
  showstringspaces=false,
  upquote=true,
  columns=fullflexible
}

% ---------- Helper macros ----------
\newcommand{\codefile}[1]{\texttt{#1}}
\newcommand{\kbd}[1]{\texttt{#1}}

% ---------- Code frame environment ----------
% Automatski dodaje fragile i shrink=5, koristi se umjesto \begin{frame}
% Primjer: \begin{codeframe}{Naslov} ... \end{codeframe}
% Primjer s opcijama: \begin{codeframe}[shrink=15]{Naslov} ... \end{codeframe}
\newenvironment{codeframe}[2][]
  {\begin{frame}[fragile,shrink=5,#1]{#2}}
  {\end{frame}}

% ============================================================

\documentclass[aspectratio=149,11pt]{beamer}
\usetheme{CambridgeUS}

\usepackage{fontspec}              % radi samo sa XeLaTeX/LuaLaTeX
\setsansfont{Latin Modern Sans}    % sans-serif font
\setmonofont{Latin Modern Mono}    % monospace font za kod
\usepackage[croatian]{babel}
\usepackage{csquotes}              % \enquote{} → hrvatski navodnici „..."
\usepackage{lmodern}
\usepackage{booktabs}
\usepackage{listings}
\usepackage{xcolor}
\usepackage{hyperref}

% ---------- Autor i logo ----------
\author{mag.ing. Josip Torić (SIT UNIZD)}
\titlegraphic{\includegraphics[height=2.5cm]{../common/SIT_logo_tekst1.pdf}}

% ---------- Listings ----------
\definecolor{codebg}{RGB}{248,248,248}
\definecolor{codeframeclr}{RGB}{220,220,220}
\lstdefinelanguage{bash}{
  morekeywords={sudo,apt,brew,pip,pipx,uv,winget,python,git,flask,source,deactivate,code,gunicorn,docker,compose,heroku,render,railway,node,npm,pnpm,yarn,uvicorn,alembic,pytest,ruff,black,precommit},
  sensitive=false,
  morecomment=[l]{\#}
}
\lstdefinelanguage{python}{
  morekeywords={from,import,def,class,return,if,elif,else,for,in,while,try,except,raise,with,as,not,and,or,is,lambda,yield,True,False,None,await,async},
  sensitive=true,
  morecomment=[l]{\#},
  morestring=[b]',
  morestring=[b]"
}
\lstdefinelanguage{sql}{
  morekeywords={SELECT,FROM,WHERE,INSERT,INTO,VALUES,UPDATE,SET,DELETE,CREATE,TABLE,PRIMARY,KEY,FOREIGN,REFERENCES,UNIQUE,NOT,NULL,INDEX,ON,ALTER,DROP,CONSTRAINT},
  sensitive=false,
  morecomment=[l]{--}
}
\lstdefinelanguage{json}{
  morecomment=[l]{//},
  morestring=[b]",
  showstringspaces=false
}
\lstdefinelanguage{typescript}{
  morekeywords={import,export,from,const,let,var,function,return,if,else,for,of,in,while,switch,case,break,async,await,try,catch,finally,throw,new,class,this,extends,implements,interface,type,enum,readonly,public,private,protected,static,get,set,true,false,null,undefined,void,string,number,boolean,any,unknown,never,typeof,keyof,as},
  sensitive=true,
  morecomment=[l]{//},
  morecomment=[s]{/*}{*/},
  morestring=[b]',
  morestring=[b]",
  morestring=[b]`
}
\lstalias{javascript}{typescript}
\lstset{
  backgroundcolor=\color{codebg},
  basicstyle=\ttfamily\footnotesize,
  frame=single,
  rulecolor=\color{codeframeclr},
  frameround=tttt,
  breaklines=true,
  tabsize=2,
  showstringspaces=false,
  upquote=true,
  columns=fullflexible
}

% ---------- Helper macros ----------
\newcommand{\codefile}[1]{\texttt{#1}}
\newcommand{\kbd}[1]{\texttt{#1}}

% ---------- Code frame environment ----------
% Automatski dodaje fragile i shrink=5, koristi se umjesto \begin{frame}
% Primjer: \begin{codeframe}{Naslov} ... \end{codeframe}
% Primjer s opcijama: \begin{codeframe}[shrink=15]{Naslov} ... \end{codeframe}
\newenvironment{codeframe}[2][]
  {\begin{frame}[fragile,shrink=5,#1]{#2}}
  {\end{frame}}


\title[Razvoj web aplikacija]{Razvoj web aplikacija - Predavanje 3 \\ Autentikacija i sigurnost (JWT)}
\date{19.10.2026.}

\begin{document}

\frame{\titlepage}

% ------------------------------------------------------------
\begin{frame}{Plan predavanja (1/2)}
	\begin{enumerate}
		\item Zašto autentikacija nije trivijalna (gdje greška boli)
		\item Hashiranje lozinki: bcrypt i zašto ne MD5/SHA
		\item JWT struktura: header, payload, signature
		\item Access token i refresh token: uloge i životni ciklus
		\item Login i refresh endpointi u FastAPI-ju
	\end{enumerate}
\end{frame}

% ------------------------------------------------------------
\begin{frame}{Plan predavanja (2/2)}
	\begin{enumerate}
		\setcounter{enumi}{5}
		\item Dependency \texttt{get\_current\_user}
		\item Token storage: localStorage vs HttpOnly cookie
		\item XSS i CSRF: što su i koji storage što štiti
		\item Testovi za auth endpointe
	\end{enumerate}

	\vspace{0.6em}
	Na kraju predavanja imate funkcionalan login flow: korisnik se može prijaviti, dobiti token i koristiti zaštićene endpointe.
\end{frame}

% ------------------------------------------------------------
\begin{frame}{Ishodi učenja}
	Nakon ovog predavanja student može:
	\begin{itemize}
		\item objasniti zašto lozinke hashiramo i što je salt
		\item opisati strukturu JWT tokena i što svaki dio radi
		\item objasniti razliku između access i refresh tokena
		\item implementirati \texttt{get\_current\_user} dependency u FastAPI-ju
		\item usporediti localStorage i HttpOnly cookie kao strategije pohrane
		\item objasniti XSS i CSRF napade i koji storage što štiti
		\item napisati test koji provjerava login i zaštićeni endpoint
	\end{itemize}
\end{frame}

% ------------------------------------------------------------
\begin{frame}{Kontekst: gdje smo u projektu?}
	\begin{itemize}
		\item Predavanje 2: imamo tablice \textbf{Club} i \textbf{User}, migracije i seed
		\item Seed je kreirao admin korisnika i demo klub korisnika
		\item Danas: ti korisnici se mogu \textbf{prijaviti} i dobiti token
	\end{itemize}

	\vspace{0.6em}
	Sljedeći korak (Predavanje 4): autorizacija - \textbf{tko smije što raditi} (role i ownership).
	Autentikacija je samo: \textbf{tko si ti}.
\end{frame}

% ------------------------------------------------------------
\begin{frame}{Zašto autentikacija nije trivijalna?}
	\begin{itemize}
		\item Greška u auth = napadač dobiva pristup tuđim podacima
		\item Najčešće greške iz prakse:
		\begin{itemize}
			\item lozinka pohranjena u plaintext ili slabo hashirana
			\item JWT tajni ključ je \enquote{secret} ili prazan string
			\item token se nikad ne invalidira (nema expiry)
			\item osjetljivi podaci u JWT payloadu - payload je čitljiv svakome
		\end{itemize}
	\end{itemize}

	\vspace{0.6em}
	Auth je jedan od rijetkih dijelova gdje \enquote{radi mi} nije dovoljno. Mora raditi i \textbf{sigurno}.
\end{frame}

% ------------------------------------------------------------
\begin{frame}{Lozinke: zašto ne smijemo pohranjivati plaintext?}
	\begin{itemize}
		\item Ako napadač ukrade bazu: odmah ima pristup svim računima
		\item Isti problem vrijedi za slabe hasheve (MD5, SHA-1, SHA-256 bez salta)
		\begin{itemize}
			\item rainbow tablice i brute force su trivijalni
		\end{itemize}
		\item Rješenje: \textbf{spori hash s per-password saltom}
	\end{itemize}

	\vspace{0.6em}
	Baza se može ukrasti. Lozinke moraju biti beskorisne napadaču čak i tada.
\end{frame}

% ------------------------------------------------------------
\begin{frame}{bcrypt: kako radi (intuitivno)}
	\begin{itemize}
		\item bcrypt generira \textbf{random salt} pri svakom hashiranju
		\item Salt se ugradi u hash string - nema potrebe za posebnim stupcem
		\item \textbf{Work factor} (cost): koliko je hash računski skup
		\begin{itemize}
			\item viši cost = sporiji napad, ali i sporija provjera
			\item preporučeni cost: 12 - dobar balans sigurnosti i brzine
		\end{itemize}
		\item Ista lozinka $\rightarrow$ različit hash svaki put (salt je drugačiji)
	\end{itemize}

	\vspace{0.6em}
	Na obrani pitam: što je salt i zašto se zbog njega isti hash nikad ne ponavlja za istu lozinku?
\end{frame}

% ------------------------------------------------------------
\begin{frame}[fragile,shrink=5]{bcrypt u Pythonu}
\begin{lstlisting}[language=python]
import bcrypt

def hash_password(plain: str) -> str:
    return bcrypt.hashpw(
        plain.encode(), bcrypt.gensalt()
    ).decode()

def verify_password(plain: str, hashed: str) -> bool:
    return bcrypt.checkpw(
        plain.encode(), hashed.encode()
    )
\end{lstlisting}

	\vspace{0.4em}
	Koristimo bcrypt direktno, ne passlib wrapper - passlib ima problem kompatibilnosti s bcrypt $\geq$ 4.1.
	\texttt{checkpw} interno radi constant-time usporedbu - to je zaštita od timing napada.
\end{frame}

% ------------------------------------------------------------
\begin{frame}{JWT: što je i što nije}
	\begin{itemize}
		\item JWT = JSON Web Token - standardizirani format (RFC 7519)
		\item Token je \textbf{potpisan}, ali \textbf{nije enkriptiran} (osim JWE)
		\begin{itemize}
			\item payload se može pročitati bez ključa (\texttt{base64url decode})
			\item stoga: ne stavljajte tajne podatke u payload
		\end{itemize}
		\item Potpis garantira: \textbf{token nije modificiran} i \textbf{izdao ga je naš server}
	\end{itemize}

	\vspace{0.6em}
	Česta greška: misleći da je JWT \enquote{enkriptiran}, developer u payload stavlja lozinku ili osjetljive osobne podatke.
\end{frame}

% ------------------------------------------------------------
\begin{frame}{JWT struktura: header.payload.signature}
	Token izgleda ovako (tri dijela odvojena točkom):
	\vspace{0.3em}
	\begin{center}
		\texttt{\textcolor{red}{xxxxx}.\textcolor{blue}{yyyyy}.\textcolor{violet}{zzzzz}}
	\end{center}

	\vspace{0.6em}
	\begin{itemize}
		\item \textcolor{red}{\textbf{Header}}: algoritam i tip tokena
		\begin{itemize}
			\item \texttt{\{ "alg": "HS256", "typ": "JWT" \}}
		\end{itemize}
		\item \textcolor{blue}{\textbf{Payload}}: claims (podaci o korisniku i tokenu)
		\begin{itemize}
			\item \texttt{sub} (subject), \texttt{exp} (expiry), \texttt{iat} (issued at), custom claims
		\end{itemize}
		\item \textcolor{violet}{\textbf{Signature}}: \texttt{HMAC-SHA256(base64(header) + "." + base64(payload), secret)}
	\end{itemize}
\end{frame}

% ------------------------------------------------------------
\begin{frame}{JWT Claims: što stavljamo u payload?}
	\begin{itemize}
		\item \textbf{Standardni claims} (preporučeni):
		\begin{itemize}
			\item \texttt{sub}: identifikator korisnika (npr. \texttt{"42"})
			\item \texttt{exp}: Unix timestamp isteka tokena
			\item \texttt{iat}: Unix timestamp izdavanja
			\item \texttt{iss}: issuer (npr. \texttt{"sit-unizd"})
		\end{itemize}
		\item \textbf{Custom claims} (naše aplikacije):
		\begin{itemize}
			\item \texttt{role}: \texttt{"admin"} ili \texttt{"club"}
			\item \texttt{club\_id}: FK za ownership provjere
		\end{itemize}
	\end{itemize}

	\vspace{0.6em}
	Payload mora biti \textbf{minimalan}. Samo ono što treba za autorizaciju - ništa više.
\end{frame}

% ------------------------------------------------------------
\begin{frame}{Access token vs Refresh token}
	\begin{tabular}{@{}lll@{}}
		\toprule
		& \textbf{Access token} & \textbf{Refresh token} \\
		\midrule
		Trajanje        & kratko (15 min)       & dugo (7-30 dana) \\
		Gdje se šalje   & svaki API request     & samo na \texttt{/auth/refresh} \\
		Gdje se čuva    & memory / localStorage & localStorage / HttpOnly cookie \\
		Ako iscuri      & kratak prozor rizika  & ozbiljan problem \\
		\bottomrule
	\end{tabular}

	\vspace{0.6em}
	Access token je \textbf{kratkotrajan} jer ga ne možemo invalidirati bez centralnog stanja.
	Refresh token omogućava \enquote{tiho} obnavljanje bez ponovne prijave.
\end{frame}

% ------------------------------------------------------------
\begin{frame}{Životni ciklus tokena (flow)}
	\begin{enumerate}
		\item Korisnik šalje username + lozinku na \texttt{POST /auth/login}
		\item Backend verificira lozinku, izdaje \textbf{access} i \textbf{refresh} token
		\item Frontend pohranjuje tokene - o tome u nastavku
		\item Na svakom zaštićenom requestu: šalje access token u \texttt{Authorization: Bearer ...}
		\item Kad access token istekne (401): frontend tiho poziva \texttt{POST /auth/refresh}
		\item Backend verificira refresh token, izdaje novi access token
		\item Korisnik ne primjećuje ništa
	\end{enumerate}
\end{frame}

% ------------------------------------------------------------
\begin{frame}[fragile, shrink=15]{JWT utils: izdavanje i verifikacija}
\begin{lstlisting}[language=python]
from datetime import datetime, timedelta, timezone
import jwt

def create_access_token(user_id: int, role: str,
                        club_id: int | None = None) -> str:
    payload = {
        "sub": str(user_id),
        "role": role,
        "club_id": club_id,
        "exp": datetime.now(timezone.utc) + timedelta(minutes=15),
        "iat": datetime.now(timezone.utc),
        "iss": settings.JWT_ISSUER,
    }
    return jwt.encode(payload, settings.JWT_SECRET, algorithm="HS256")

def decode_token(token: str) -> dict:
    try:
        return jwt.decode(token, settings.JWT_SECRET,
                          algorithms=["HS256"],
                          issuer=settings.JWT_ISSUER,
                          options={"require": ["exp"]})
    except jwt.ExpiredSignatureError:
        raise UnauthorizedError("Token expired")
    except jwt.InvalidTokenError:
        raise UnauthorizedError("Invalid token")
\end{lstlisting}

	\vspace{0.4em}
	Koristimo \textbf{PyJWT} - aktivno održavana biblioteka. Razlikovanje \texttt{ExpiredSignatureError} od \texttt{InvalidTokenError} omogućuje da frontend zna treba li refresh (istekao) ili ponovni login (nevažeći).
\end{frame}

% ------------------------------------------------------------
\begin{frame}[fragile, shrink=10]{Login endpoint}
\begin{lstlisting}[language=python]
from fastapi import APIRouter, Depends
from sqlalchemy.ext.asyncio import AsyncSession
from app.core.deps import get_db
from app.schemas.auth import LoginRequest, TokenResponse
from app.services.auth import authenticate_user, create_tokens

router = APIRouter()

@router.post("/login", response_model=TokenResponse)
async def login(body: LoginRequest,
                db: AsyncSession = Depends(get_db)):
    user = await authenticate_user(db, body.username, body.password)
    access_token, refresh_token = create_tokens(user)
    return TokenResponse(access_token=access_token,
                         refresh_token=refresh_token)
\end{lstlisting}

	\vspace{0.4em}
	\texttt{authenticate\_user} je u service layeru, ne u routeru.
\end{frame}

% ------------------------------------------------------------
\begin{frame}[fragile,shrink=5]{Refresh endpoint}
\begin{lstlisting}[language=python]
@router.post("/refresh", response_model=TokenResponse)
async def refresh(body: RefreshRequest,
                  db: AsyncSession = Depends(get_db)):
    payload = decode_token(body.refresh_token)
    user = await get_user_by_username(db, payload["sub"])
    if not user or not user.is_active:
        raise UnauthorizedError("User not found or inactive")
    access_token, new_refresh = create_tokens(user)
    return TokenResponse(access_token=access_token,
                         refresh_token=new_refresh)
\end{lstlisting}

	\vspace{0.4em}
	Razmislite do obrane: zašto provjeravamo \texttt{is\_active}, a ne samo potpis tokena?
\end{frame}

% ------------------------------------------------------------
\begin{frame}[fragile,shrink=5]{Dependency: \texttt{get\_current\_user}}
\begin{lstlisting}[language=python]
from fastapi import Depends, Header
from app.core.errors import UnauthorizedError

async def get_current_user(
    authorization: str = Header(...),
    db: AsyncSession = Depends(get_db)
) -> User:
    if not authorization.startswith("Bearer "):
        raise UnauthorizedError("Missing token")
    token = authorization.removeprefix("Bearer ")
    payload = decode_token(token)
    user = await get_user_by_id(db, payload["sub"])
    if not user or not user.is_active:
        raise UnauthorizedError("User not found")
    return user
\end{lstlisting}
\end{frame}

% ------------------------------------------------------------
\begin{frame}[fragile,shrink=5]{Korištenje dependencyja u routeru}
\begin{lstlisting}[language=python]
@router.get("/me")
async def get_me(current_user: User = Depends(get_current_user)):
    return {"id": current_user.id, "username": current_user.username,
            "role": current_user.role}

@router.get("/clubs")
async def list_clubs(
    db: AsyncSession = Depends(get_db),
    current_user: User = Depends(get_current_user)
):
    # svaki korisnik koji je prijavljen smije vidjeti klubove
    return await get_all_clubs(db)
\end{lstlisting}

	\vspace{0.4em}
	Provjera \textbf{role i ownership} dolazi u Predavanju 4.
\end{frame}

% ------------------------------------------------------------
\begin{frame}{Zašto dependency, a ne middleware za auth?}
	\begin{itemize}
		\item Middleware se izvršava za \textbf{sve} requestove (i \texttt{/health}, i \texttt{/docs})
		\item Dependency se izvršava samo gdje ga \textbf{deklariramo}
		\item Dependency je \textbf{testabilniji}: lako ga možemo zamijeniti u testovima
		\item FastAPI dependency injection je eksplicitan:
		\begin{itemize}
			\item odmah vidimo koji endpoint treba autentikaciju
		\end{itemize}
	\end{itemize}

	\vspace{0.6em}
	Pitanje za razmišljanje: kako bi testirao endpoint koji koristi \texttt{get\_current\_user}? Odgovor stiže u Predavanju 11.
\end{frame}

% ------------------------------------------------------------
\begin{frame}{Token storage: gdje čuvamo tokene na frontendu?}
	Dvije uobičajene opcije:

	\vspace{0.6em}
	\begin{tabular}{@{}lll@{}}
		\toprule
		& \textbf{localStorage} & \textbf{HttpOnly cookie} \\
		\midrule
		Pristup JS-u         & da                  & ne \\
		XSS ranjivost        & visoka              & niska \\
		CSRF ranjivost       & niska               & umjerena \\
		Implementacija       & jednostavna         & složenija \\
		Pogodnost za SPA     & visoka              & zahtijeva \texttt{credentials} \\
		\bottomrule
	\end{tabular}

	\vspace{0.6em}
	U projektu koristimo \textbf{localStorage} za oba tokena - jednostavnija implementacija za SPA. HttpOnly cookie je sigurnija produkcijska alternativa, pokazat ćemo i nju.
\end{frame}

% ------------------------------------------------------------
\begin{frame}{XSS (Cross-Site Scripting): što je i zašto boli}
	\begin{itemize}
		\item Napadač ubaci zloćudni JavaScript u stranicu (npr. kroz komentar, URL)
		\item Taj JS se izvršava u browseru žrtve - ima isti pristup kao vaš kod
		\item Ako je token u \texttt{localStorage}:
		\begin{itemize}
			\item \texttt{document.cookie} nema token, ali \texttt{localStorage.getItem()} ga ima
			\item napadač krade token s jednom linijom koda
		\end{itemize}
		\item Zaštita: HttpOnly cookie (JS ga ne može pročitati), CSP header, sanitizacija outputa
	\end{itemize}
\end{frame}

% ------------------------------------------------------------
\begin{frame}{CSRF (Cross-Site Request Forgery): što je i zašto boli}
	\begin{itemize}
		\item Napadač nagovori korisnika da posjeti stranicu koja šalje request \textbf{u korisnikovo ime}
		\item Browser automatski šalje cookieje uz request (uključujući auth cookie)
		\item Ako koristite HttpOnly cookie bez CSRF zaštite: ranjivi ste na CSRF
		\item Zaštita: \texttt{SameSite=Strict/Lax} na cookie, CSRF token, provjera \texttt{Origin} headera
	\end{itemize}

	\vspace{0.6em}
	localStorage + XSS = krađa tokena. \\
	HttpOnly cookie + CSRF = izvršavanje akcija. \\
	Nema savršene opcije - bira se manji rizik.
\end{frame}

% ------------------------------------------------------------
\begin{frame}{Praktična odluka za naš projekt}
	\begin{itemize}
		\item \textbf{Access i refresh token}: \texttt{localStorage}
		\begin{itemize}
			\item preživljava refresh stranice - korisnik ostaje prijavljen
			\item frontend implementaciju radimo u Predavanju 8
		\end{itemize}
		\item Cijena: token je dostupan JavaScriptu, XSS ga može ukrasti
		\item Zato: kratak access token (15 min), Vue auto-escaping, nikad \texttt{v-html} s korisničkim sadržajem
	\end{itemize}

	\vspace{0.6em}
	U produkcijskim sustavima refresh token ide u \textbf{HttpOnly cookie} - JS ga ne može pročitati. Za naš projekt localStorage je svjesni trade-off, ne neznanje.

	\vspace{0.4em}
	Na obrani pitam: što napadač dobiva XSS-om ako su tokeni u localStorageu? Kako HttpOnly cookie mijenja rizik?
\end{frame}

% ------------------------------------------------------------
\begin{frame}[fragile,shrink=5]{Produkcijska alternativa: HttpOnly cookie (FastAPI)}
\begin{lstlisting}[language=python]
from fastapi import Response

@router.post("/login")
async def login(body: LoginRequest, response: Response,
                db: AsyncSession = Depends(get_db)):
    user = await authenticate_user(db, body.username, body.password)
    access_token, refresh_token = create_tokens(user)
    response.set_cookie(
        key="refresh_token",
        value=refresh_token,
        httponly=True,
        secure=True,       # samo HTTPS u produkciji
        samesite="lax",
        max_age=60 * 60 * 24 * 30,  # 30 dana
    )
    return {"access_token": access_token}
\end{lstlisting}
\end{frame}

% ------------------------------------------------------------
\begin{frame}{Testovi za auth: što pokrivamo?}
	\begin{itemize}
		\item \textbf{Uspješan login}: točni kredencijali $\rightarrow$ 200 + token
		\item \textbf{Krivi podaci}: pogrešna lozinka $\rightarrow$ 401
		\item \textbf{Nepostojeći username}: $\rightarrow$ 401 (isti error, ne otkrivamo postoji li)
		\item \textbf{Zaštićeni endpoint bez tokena}: $\rightarrow$ 401
		\item \textbf{Zaštićeni endpoint s tokenom}: $\rightarrow$ 200
		\item \textbf{Istekli token}: $\rightarrow$ 401
		\item \textbf{Refresh}: valjan refresh token $\rightarrow$ novi access token
	\end{itemize}

	\vspace{0.6em}
	Zašto isti error za krivu lozinku i nepostojeći username? Enumeracija korisnika: napadač bi mogao otkriti koji usernamei postoje u sustavu.
\end{frame}

% ------------------------------------------------------------
\begin{frame}[fragile, shrink=10]{Primjer testa (pytest + TestClient)}
\begin{lstlisting}[language=python]
from fastapi.testclient import TestClient
from app.main import app

client = TestClient(app)

def test_login_success(db_with_seed):
    r = client.post("/auth/login",
                    json={"username": "admin",
                          "password": "admin123"})
    assert r.status_code == 200
    assert "access_token" in r.json()

def test_login_wrong_password(db_with_seed):
    r = client.post("/auth/login",
                    json={"username": "admin",
                          "password": "wrong"})
    assert r.status_code == 401

def test_protected_without_token():
    r = client.get("/auth/me")
    assert r.status_code == 401
\end{lstlisting}
\end{frame}

% ------------------------------------------------------------
\begin{frame}[fragile,shrink=5]{Override dependencyja u testovima}
\begin{lstlisting}[language=python]
from app.core.deps import get_current_user
from app.models.user import User

def fake_admin():
    return User(id=1, username="admin",
                role="admin", is_active=True)

app.dependency_overrides[get_current_user] = fake_admin

def test_me_endpoint():
    r = client.get("/auth/me")
    assert r.status_code == 200
    assert r.json()["role"] == "admin"

# uvijek očistite override nakon testa
app.dependency_overrides = {}
\end{lstlisting}

	\vspace{0.4em}
	Testiramo endpoint \textbf{neovisno} o auth logici.
\end{frame}

% ------------------------------------------------------------
\begin{frame}{Sigurnost JWT tajnog ključa}
	\begin{itemize}
		\item Tajni ključ \textbf{mora biti} nasumičan i dugačak (min. 32 bajta)
		\begin{itemize}
			\item \texttt{openssl rand -hex 32}
		\end{itemize}
		\item Nikad ne smije biti u kodu ili u repozitoriju
		\item Env var: \texttt{JWT\_SECRET} iz \codefile{.env} (koji je u \codefile{.gitignore})
		\item Ako se kompromitira: \textbf{svi tokeni su nevaljani} - morate rotirati ključ
	\end{itemize}

	\vspace{0.6em}
	Zadatak za znatiželjne: kako bi invalidirao sve aktivne tokene odjednom, a da ne mijenjaš \texttt{JWT\_SECRET}?
\end{frame}

% ------------------------------------------------------------
\begin{frame}{Brute force na login: rate limiting}
	Login endpoint je meta za pogađanje lozinki. bcrypt usporava svaki pojedinačni pokušaj, ali ne sprječava napadača da ih pošalje tisuće.

	\vspace{0.6em}
	\begin{tabular}{@{}ll@{}}
		\toprule
		\textbf{Mjera} & \textbf{Što sprječava} \\
		\midrule
		bcrypt s visokim costom & offline crackanje ukradenih hasheva \\
		rate limiting & online pogađanje kroz login endpoint \\
		\bottomrule
	\end{tabular}

	\vspace{0.6em}
	Rate limiting ograničava broj pokušaja po IP-u ili usernameu, npr. 5 pokušaja u minuti. U FastAPI ekosustavu to radi biblioteka \texttt{slowapi}. Za studentski projekt nije obavezan, ali na obrani morate znati objasniti zašto bi ga produkcijski sustav imao.
\end{frame}

% ------------------------------------------------------------
\begin{frame}{Token invalidacija: ograničenja stateless JWT-a}
	\begin{itemize}
		\item JWT je \textbf{stateless}: server ne prati koji su tokeni aktivni
		\item Ako se korisnik odjavi ili promijeni lozinku, stari token i dalje vrijedi do expiry-ja
		\item Opcije za invalidaciju:
		\begin{itemize}
			\item \textbf{Kratki expiry} (access token): prozor rizika je mali
			\item \textbf{Revocation list} (blacklist): centralna lista invalidnih tokena (Redis)
			\item \textbf{Token version} u bazi: ako se ne poklapa, token je nevaljan
		\end{itemize}
	\end{itemize}

	\vspace{0.6em}
	Za ovaj projekt: kratki access token + refresh token je dovoljan. Revocation list je \enquote{nice to have}.
\end{frame}

% ------------------------------------------------------------
\begin{frame}{Commitovi za Predavanje 3}
\begin{itemize}
  \item \texttt{feat(api): add bcrypt helpers, auth schemas and user repository}
  \item \texttt{feat(api): add JWT utilities (create/decode access + refresh tokens)}
  \item \texttt{feat(api): add auth service and login/refresh endpoints}
  \item \texttt{feat(api): add get\_current\_user dependency and /auth/me endpoint}
  \item \texttt{test(api): add auth tests (login, refresh, expired token, protected endpoint)}
\end{itemize}

\vspace{0.6em}
\textbf{Na satu:} prolazimo diffeve i testiramo sa Swaggerom - login, token, zaštićeni endpoint.
\end{frame}

% ------------------------------------------------------------
\begin{frame}{Što ide sljedeće (Predavanje 4)}
	\begin{itemize}
		\item Autorizacija: tko smije što raditi
		\item Role check: admin vs club korisnik
		\item Ownership: korisnik smije mijenjati samo svoje
		\item Dependency \texttt{require\_role} i \texttt{require\_admin}
		\item Greške: 401 vs 403 - razlika i kad što vraćamo
	\end{itemize}
\end{frame}

\end{document}
